\documentclass[11pt]{article}

\usepackage[margin=1.1in]{geometry}
\usepackage{amsmath,amssymb}
\usepackage{booktabs}
\usepackage{array}
\usepackage{enumitem}
\usepackage{microtype}
\usepackage[authoryear,round]{natbib}
\usepackage[hidelinks]{hyperref}
\usepackage{cleveref}

\title{The Mind of God:\\
Algorithmic Agnosticism and Religious Belief}
\author{Christopher S. Taylor}
\date{Version 1.0 --- preprint, not peer-reviewed --- November 2026\\[0.5ex] \small Licensed under CC BY 4.0 --- \url{https://github.com/extrafox/the-mind-of-god}}

\begin{document}
\maketitle

\begin{abstract}
Hawking hoped that a complete theory would let us know the mind of God. \emph{Algorithmic
Agnosticism} argued that no theory of everything can be known to be final. This paper asks the
parallel question of religious belief: not whether God exists, but whether any worldview can be
established as the true account of what is. It adds one premise, in two clauses: a worldview that
claims to be true is an account in the sense of \emph{Algorithmic Agnosticism}; and any account
that goes beyond the evidence, a worldview included, has a rival consistent with the evidence and
the truths of reason on which it is false. A second premise classifies every element of a classic
argument as evidence, an account, a judgment, an imposed constraint or a truth of reason. With it,
the paper locates the decisive element of twenty-two classic arguments for and against theism and
naturalism, and for the skeptical hypothesis. Three conclusions follow. No worldview can be
established as the true account. Worldviews that agree on all possible evidence are never
separated, even by final evidence. What selects among the worldviews the evidence leaves open is a
choice the evidence does not determine. These conclusions concern what can be established, not
what it is reasonable to believe. They amount to agnosticism in Huxley's sense: theism, naturalism
or the skeptical hypothesis may be true, but no argument from experience or reason can establish
that any of them is.
\end{abstract}

\section{Introduction}\label{sec:intro}

Stephen Hawking's \emph{A Brief History of Time} ends with a hope. If physics found a complete
theory and we could then understand why the universe exists, ``it would be the ultimate triumph
of human reason---for then we would know the mind of God'' \citep[ch.~11]{hawking1988}.
\emph{Algorithmic Agnosticism} argued that no account of everything can be known to be final, and that what
selects among the accounts the evidence leaves open is a choice \citep[§§3, 9]{taylor2026aa}.

This paper asks the parallel question of religious belief. It does not ask whether God exists. It
asks whether any worldview---theism, naturalism, the skeptical hypothesis, or the teaching of any
religious tradition---can be established as the true account of what is.

Its position has three parts. A worldview is a claim about what ultimately is, and an account
consistent with everything anyone has experienced makes it false. To establish a worldview would be
to show that the claim holds in every account the evidence leaves open. No one can do that, at any
time, and what selects among the worldviews the evidence leaves open is a choice. Stated for the worldviews at issue:

\begin{quote}
No argument from experience or reason establishes that any worldview, whether theism,
naturalism or the skeptical hypothesis, is the true account of what is. Any of them may be true,
but no one can establish that it is.
\end{quote}

The claim is one of parity. It is not that God does not exist, or that nature is all there is.

The parity matters because it cuts both ways. It removes demonstration from the theist and the
atheist alike. The proofs of natural theology and the disproofs offered against them fall under the
same verdict, so neither side can claim that reason has decided the question in its favor. It
leaves belief and faith untouched, and with them the question of what it is reasonable to believe;
the paper is silent on these. What the parity rules out is any argument's claim to settle the question for
everyone. An argument may persuade those who accept its premises and its standard of preference.
It does not settle the question for anyone who shares the same evidence and declines them.

The core claim has neighbors. \citet{oppy2006} argues that no argument about God is successful, by a
dialectical standard of success. \citet{mizrahi2017} argues that God's existence can be
established neither \emph{a priori} nor \emph{a posteriori}, a position he calls skeptheism.
\citet{alston1986} argues that the reliability of a doxastic practice cannot be shown without
epistemic circularity. \citet{hick1989} and \citet{mckim2001} argue that the universe is
religiously ambiguous. This paper differs in three ways. It derives its conclusions from the
premise and results of \emph{Algorithmic Agnosticism}, with a premise that extends its first result to worldviews. Its classification locates the
decisive element of the classic arguments in \cref{sec:arguments}. And its parity includes
naturalism and the skeptical hypothesis.

\Cref{sec:aa} restates the results of \emph{Algorithmic Agnosticism}. \Cref{sec:classification} gives the
definitions and the premises. \Cref{sec:arguments} takes up the classic arguments.
\Cref{sec:result} derives the conclusions. \Cref{sec:objections} answers objections, and
\cref{sec:conclusion} concludes.

\section{The results of \emph{Algorithmic Agnosticism}}\label{sec:aa}

This section restates the parts of \emph{Algorithmic Agnosticism} (AA) that this paper uses, in the language of
philosophy of religion. The proofs are there.

\medskip\noindent\textbf{AA-D. Classes.}\phantomsection\label{item:AAD}
\citep[§2]{taylor2026aa}
\begin{itemize}
  \item \emph{Evidence} is what a person has observed, experienced, remembered or written down.
        It is non-factive about the world: ``it seemed that God spoke'', not ``God spoke''.
  \item An \emph{account} is a statement of how the world is.
  \item A \emph{judgment} is a way of passing from evidence to an account.
  \item An \emph{imposed constraint} is a criterion for choosing among accounts that the
        evidence does not determine.
  \item \emph{Truths of reason} (this paper's term) are logic, mathematics, and conceptual
        analysis.
\end{itemize}

\medskip\noindent\textbf{AA-P1. Finite evidence.}\phantomsection\label{item:AAP1}
\citep[§2]{taylor2026aa} Every person's evidence is finite.

\medskip\noindent\textbf{AA-R1. No establishment.}\phantomsection\label{item:AAR1}
\citep[§3]{taylor2026aa}
\begin{itemize}
  \item Every finite body of evidence fits infinitely many accounts that disagree about what has
        not yet been observed.
  \item No judgment of any kind establishes one of them as the true account: any judgment that selects one as
        the true account is wrong in some account consistent with the evidence.
\end{itemize}

\medskip\noindent\textbf{AA-R2. Permanent underdetermination.}\phantomsection\label{item:AAR2}
\citep[§6]{taylor2026aa} Accounts that agree on all possible evidence are never separated by any
judgment, even by an ideal judge with all the evidence there could ever be.

\medskip\noindent\textbf{AA-R3. What remains is choice.}\phantomsection\label{item:AAR3}
\citep[§§4--5, 9]{taylor2026aa} Among the accounts that the evidence leaves open, selection is
made by an imposed constraint, and the evidence does not determine which.

\section{Definitions and premises}\label{sec:classification}

\medskip\noindent\textbf{D1. Worldview.}\phantomsection\label{item:D1}
(Definition.) A claim about what ultimately is. Every worldview goes beyond any finite body of
evidence. Examples are theism, naturalism, the skeptical hypothesis (that we are deceived, or
simulated), and the teaching of any religious tradition.

\medskip\noindent\textbf{D2. Classic argument.}\phantomsection\label{item:D2}
(Definition.) An argument offered to establish, or to refute, a worldview.

\medskip\noindent\textbf{D3. Establish.}\phantomsection\label{item:D3}
(Definition.) An argument establishes a claim from the evidence if and only if the claim is true in every
account consistent with the evidence and with the truths of reason. This is \emph{Algorithmic Agnosticism}'s
standard for a theory known to be final, applied to worldviews: the claim holds in every account the evidence leaves open.

Natural theology claims demonstration. Aquinas holds that the existence of God, though not
self-evident to us, can be demonstrated from its effects, which are known to us
\citep[I q.2 a.2]{aquinas-st}. The First Vatican Council teaches that God, the beginning and end of
all things, ``can be known with certainty'' from created things by the natural light of human
reason \citep[ch.~2]{vatican1870}. The denial has had its demonstrations too. The logical problem
of evil holds that God's existence is inconsistent with the existence of evil, and Findlay argued
that the being worship requires would have to exist necessarily, and that, if necessity belongs
only to propositions, no being can \citep{mackie1955,findlay1948}. Each side claims more than
probability. D3 names the standard that these proofs and disproofs alike set
out to meet.

\medskip\noindent\textbf{P-M. Mapping premise.}\phantomsection\label{item:PM}
(Premise.)
\begin{enumerate}[label=(\alph*)]
  \item A worldview, insofar as it claims to be true, is an account in the sense of
        \emph{Algorithmic Agnosticism}. Every route to belief, argument or not, supplies evidence, a judgment, or an
        imposed constraint.
  \item For any finite evidence, any account that goes beyond it, a worldview included, has a
        rival consistent with the evidence and the truths of reason on which it is false.
\end{enumerate}
\emph{Motivation.} Clause (a) follows the treatment of religious belief as a ``religious
hypothesis'' \citep[§X]{james1896}, or as an explanatory hypothesis assessed by its fit with the
evidence \citep{swinburne2004}. \emph{Algorithmic Agnosticism} proves for its own accounts what (b) asserts of
worldviews (AA-R1), with rivals that differ on what has not yet been observed. Those accounts do
not speak of God or nature, so for worldviews (b) is a premise.
Two rivals illustrate (b). For theism, one rival is an account on which the same course of
experience, religious experience included, arises with nothing beyond nature: the sense of God's
presence, the order of the world and the testimony of the traditions are what natural processes
would bring about. For naturalism, one rival is an account on which God creates and sustains
exactly the world the sciences describe: every law holds and every experience runs as it does,
but the world depends for its existence on God. Each rival agrees with all the evidence that the
view it opposes cites. Each is the kind of rival the classic debates already treat as live: theists answer
natural explanations of religious experience by argument, and naturalists answer a God who works
through the order of nature by argument, rather than dismissing either as incoherent. What (b)
asserts of every worldview is thus what each side already treats as live in the other's: that it
has such a rival.
\emph{Refuted by:} for (a), a worldview that claims to be true yet is no account of what is, or a
route to belief that delivers something else; for (b), an account that goes beyond some finite evidence, a worldview included, yet is
true in every account consistent with that evidence and the truths of reason. A sound ontological
argument would show this for theism, and a sound logical problem of evil for its denial.

\medskip\noindent\textbf{T-Class. Classification premise.}\phantomsection\label{item:TClass}
(Premise; used for coverage in \cref{sec:arguments}.) Every element of a classic argument is one
of five things: evidence, an account, a judgment, an imposed constraint, or a truth of reason.
\begin{itemize}
  \item \emph{Propositions} \citep[§4]{hume1748}.
        A matter of fact is either a report of experience, which is evidence, or a claim about
        what is, which is an account. A relation of ideas is a truth of reason. A claim offered
        \emph{a priori} that would exclude an account consistent with the truths of reason is not
        a truth of reason, since that account is consistent with them. It serves as a criterion for
        choosing among accounts the evidence leaves open, so it is an imposed constraint.
        Examples are the possibility of a maximally great being, the principle of
        sufficient reason, the incompatibility of God and evil, and a disputed modal axiom.
  \item \emph{Rules.} A valid deductive rule is a truth of reason. A non-deductive rule
        (induction, inference to the best explanation, probabilistic confirmation, the
        deliverance of a faculty) is a judgment. A standard for preferring one conclusion to
        another is an imposed constraint.
  \item \emph{Practical reasons}, such as those of \citet{pascal1670} and \citet{james1896},
        are reasons for choosing a belief, so they are imposed constraints.
  \item \emph{Nothing else enters an argument.} An argument consists of propositions, and of
        rules by which they are taken to support a conclusion.
\end{itemize}
\emph{Refuted by} an element of a classic argument that fits none of the five classes. Hume's fork
is itself contested, so T-Class is owned as a premise, not claimed as a truth of reason.

\section{The classic arguments}\label{sec:arguments}

This section applies T-Class (\cref{item:TClass}) case by case. For each argument it states a
standard form, names the \emph{decisive element}, the one that carries the conclusion beyond the
evidence, and gives that element's class: an account, a judgment, or an imposed constraint. By
P-M(b), a decisive premise offered \emph{a priori}, in an argument valid for its conclusion, would
exclude an account consistent with the truths of reason, so by T-Class it is an imposed constraint.
In a non-deductive argument, the rule of inference is a judgment.

\subsection{A worked example}\label{ss:worked-example}

One argument is set out here in full, to show how the classification proceeds; the others are given
more briefly. The conclusions this example previews are derived in \cref{sec:result}.

The fine-tuning argument, in its probabilistic form, has four elements
\citep{swinburne2004,collins2009}. First, the constants of nature, such as the strength of gravity
and the cosmological constant, fall in a range that permits life, and on the argument's estimates a
narrow one. Second, a life-permitting universe is far more probable if it was designed than on its
rivals: on design it is to be expected, and on a rival on which the constants have no explanation it
is a coincidence. Third, design has a prior probability that is not too low. Fourth, by Bayes'
theorem, the probability of an account given the evidence rises with how much more probable the
evidence is on that account than on its rivals, and with the account's prior probability; so, given
the first three elements, design is probable.

Each element falls in one of T-Class's classes. The first is evidence: it reports the values of the
constants as physics has measured them. The second is a judgment: it passes from the evidence to a
comparison of accounts by estimating how probable the evidence is on each, and probabilistic
confirmation is a non-deductive rule. The third is an imposed constraint: a prior ranks the accounts
before the evidence is weighed, and it is a criterion for choosing among them that the evidence does
not determine. The fourth is a valid inference, since Bayes' theorem is a theorem of probability, and
so it is a truth of reason.

Design and its rivals do not dispute the evidence, and the inference is valid, so neither of these
carries the conclusion beyond the evidence. The decisive elements are the judgment over likelihoods
and the prior.

Each is answered by a different conclusion. Against a claim that the argument establishes design,
C1 applies: an account with no designer fits the same evidence, and design is false in it. One such
account has many universes with varying constants, with observers found only in those that permit
life; another derives the constants from a deeper law. The argument as stated concludes only that
design is probable, and that conclusion C1 leaves standing. That conclusion rests on the prior as well
as the likelihoods, and for the prior C3 applies. A prior is a choice the evidence does not determine. A prior low enough, equally
consistent with the evidence, makes design improbable on the same likelihoods; one high enough makes
it probable.

C2 applies in one further case. If design and its rival agree on all possible evidence, so that
every observation that could ever be made is as probable on one as on the other, then no further
measurement separates them. If they differ somewhere on possible evidence, the case is C1's.

\subsection{Arguments for theism}\label{ss:arguments-for}

\paragraph{The kalam argument.} Whatever begins to exist has a cause; the universe began to exist;
so the universe has a cause \citep{craig1979}. Craig supports the second premise with the claim
that an actual infinite cannot exist, and argues in a further step that the cause is personal. The
decisive elements are the causal principle, the denial of an actual infinite, and the step to a
personal cause. Each is an account if offered as a generalization from experience, and an imposed
constraint if offered \emph{a priori}.

\paragraph{The contingency argument.} Every contingent fact has an explanation; the totality of
contingent facts cannot be explained by a contingent being; so there is a necessary being
\citep{pruss2009}. The decisive elements are the principle of sufficient reason and the claim that
the totality has no contingent explanation, classed as for the kalam argument.

\paragraph{The Nyāya argument.} Udayana argues that the earth, being an effect, has an intelligent
maker, Īśvara, as a pot has a potter \citep{dasti2011}. The decisive element is that the earth is an
effect requiring an intelligent maker, an account if drawn from experience and an imposed
constraint if offered \emph{a priori}.

\paragraph{The fine-tuning argument.} Its probabilistic form is set out and classified in
\cref{ss:worked-example}.

\paragraph{The moral argument.} There are objective moral facts, and a theory on which they consist
in facts about God, such as God's commands, explains them better than its nontheistic rivals
\citep{adams1987}. The decisive element is that objective moral facts exist and need a ground. It is
an account if offered as a description of how things are, and an imposed constraint if offered
\emph{a priori}.

\paragraph{The ontological argument.} In its modal form: it is possible that a maximally great being
exists; maximal greatness includes necessary existence; in the modal system S5, whatever is possibly
necessary is necessary; so a maximally great being exists \citep{plantinga1974}; compare
\citet[ch.~2--3]{anselm-proslogion}. The decisive elements are the possibility premise and the modal
axiom S5. Both are offered \emph{a priori}, and so are imposed constraints.

\paragraph{The arguments from consciousness and reason; the trademark argument.} Mind, rational
insight, or the idea of an infinite being cannot arise from unthinking nature, so they come from a
mind; Descartes argues that the idea of the infinite must have been placed in us by an infinite cause
\citep[Meditation~III]{descartes1641}. The decisive element is that mind, or the idea of the
infinite, cannot arise from nature: an account, or an imposed constraint if offered as a causal
principle held \emph{a priori}.

\paragraph{The argument from religious experience.} It seems to many people that God is present to
them, and by the principle of credulity what seems present probably is present, unless there is
reason to think otherwise \citep{swinburne2004}. Alston argues that it is practically rational to
engage in a socially established doxastic practice, the practice of forming beliefs on such
experiences no less than sense perception, and concedes that its reliability cannot be shown
without epistemic circularity \citep{alston1991}. The decisive elements are the principle of
credulity, an imposed constraint, and, for Alston, this parity principle of practical rationality,
a practical reason and so an imposed constraint.

\paragraph{Testimony and revelation.} Scripture and tradition transmit what God revealed, and the
transmission is faithful (for example, the Bible and the Qur'an; \citealp{bible,quran}). The
decisive element is that the transmission is reliable and its origin divine, an account.

\paragraph{Miracles, and Hume's argument against them.} For: God's action would raise the
probability of a reported event that nature does not explain \citep{swinburne2004}. Against: no
testimony suffices to establish a miracle unless its falsehood would be more miraculous than the
fact it reports \citep[§10.13]{hume1748}. The decisive element for the miracle is that God's action
raises the probability of the report, a judgment. The decisive element for Hume is his maxim, a
standard for preferring one conclusion to another, and so an imposed constraint.

\paragraph{First-person revelation.} The witness has the experience directly: Moses at the bush that
burned and was not consumed (Exodus 3:2--14), or Joseph Smith seeing two Personages in a pillar of
light \citep[1:16--17]{jsh1838}. The decisive element is that the experience is veridical, that it
was God who spoke, an account.

\paragraph{Inner witness and proper basicality.} A sense of the divine, or the inner testimony of the
Spirit, gives rise to belief in God directly, without inference \citep[I.3.1, I.7.4]{calvin1559},
and a belief so formed, by a faculty working as designed, may be properly basic
\citep{plantinga2000}. The decisive element is that the faculty is reliable, an account.

\paragraph{\emph{Consensus gentium}.} Belief in God is held across peoples and ages, and such
widespread agreement is offered as evidence of its truth; \citet{kelly2011} examines a modest version
of the argument. The decisive element is that convergence of belief tracks the truth, an account.

\paragraph{Pascal's Wager; James's will to believe.} Where the evidence leaves the question open, the
expected value of believing favors belief \citep{pascal1670}; or, where the option is live, forced
and momentous, one has the right to let one's passional nature decide \citep{james1896}. The
decisive element is expected value, or the right to believe, an imposed constraint (a choice).

\subsection{Arguments against theism}\label{ss:arguments-against}

\paragraph{The logical problem of evil.} A good thing eliminates evil as far as it can, and there are
no limits to what an omnipotent being can do; so an omnipotent, wholly good being would eliminate
evil completely; evil exists; so there is no such being \citep{mackie1955}. The decisive element is
that God and evil are incompatible, offered \emph{a priori}, and so an imposed constraint.

\paragraph{The evidential problem of evil.} There exist instances of intense suffering that an
omnipotent, omniscient being could have prevented without losing some greater good or permitting
some evil equally bad or worse, and such a being, if wholly good, would prevent them
\citep{rowe1979}. In a second form, the facts about pain and pleasure are more probable on the
hypothesis of indifference than on theism \citep{draper1989}. The decisive element for Rowe is that
such suffering exists, an account. For Draper it is the comparison of probabilities, a judgment,
with the prior an imposed constraint.

\paragraph{Divine hiddenness.} A perfectly loving God would ensure that everyone capable of a
relationship with him is in a position to believe that he exists, so there would be no reasonable,
or inculpable, nonbelief; there is such nonbelief; so there is no such God \citep{schellenberg1993}.
The decisive element is that perfect love entails that there is no reasonable nonbelief. It is an
account if it states what love is in fact like, and an imposed constraint if offered
\emph{a priori}.

\paragraph{Findlay's disproof.} An adequate object of worship must exist necessarily, and no
existential proposition can be necessary, so there is no such object \citep{findlay1948}. The
decisive element is that no existential proposition can be necessary, offered \emph{a priori}, and
so an imposed constraint.

\paragraph{Occam's razor.} An account of the world without God explains the same evidence with fewer
entities, so it is to be preferred (on such appeals to parsimony, see \citealp{sober2015}). The
decisive element is that the simpler account is to be preferred, a standard for preferring one
conclusion to another, and so an imposed constraint.

\subsection{Naturalism and the skeptical hypothesis}\label{ss:arguments-naturalism}

\paragraph{Causal closure.} Every physical effect has a sufficient physical cause, and Papineau
argues from this to physicalism \citep{papineau2001}; extended, it leaves nothing beyond nature that
acts in the physical world. The decisive element is that every physical effect has a sufficient
physical cause, an account.

\paragraph{The evolutionary argument against naturalism.} If naturalism and evolution are both true,
the probability that our cognitive faculties are reliable is low or inscrutable; so whoever accepts
both has a defeater for every belief, including naturalism \citep[ch.~12]{plantinga1993}. The
decisive elements are that this probability is low or inscrutable, an account, and the principle
that such a probability defeats the beliefs the faculties produce, an imposed constraint.

\paragraph{The simulation argument.} Either almost no civilizations reach the stage at which they can
run detailed simulations of minds like ours, or almost none that reach it choose to, or we are almost
certainly living in a simulation \citep{bostrom2003}. The decisive elements are the assumptions
about simulators, that minds like ours can be simulated and what mature civilizations could and
would do, an account; and the bland indifference principle, that one's credence in being simulated
should match the fraction of observers like oneself who are, an imposed constraint.

\section{The conclusions}\label{sec:result}

\noindent\textbf{C1. No worldview can be established as the true account.}\phantomsection\label{item:C1}
(From D1, D3, P-M(b) and AA-P1.)
Let W be a worldview.
\begin{itemize}
  \item W goes beyond the evidence (D1).
  \item By P-M(b), a rival R consistent with the evidence and the truths of reason makes W false.
  \item By D3, no argument establishes W, since W is false in R.
  \item This holds at every time, because evidence stays finite (AA-P1).
\end{itemize}
No premise about premises transmitting establishment is needed, since D3 concerns only the
accounts consistent with the evidence.

Stated for the worldviews at issue:

\begin{quote}
No argument from experience or reason establishes that any worldview, whether theism,
naturalism or the skeptical hypothesis, is the true account of what is. Any of them may be true,
but no one can establish that it is.
\end{quote}

This is agnosticism in Huxley's sense, ``not a creed, but a method'': ``do not pretend that
conclusions are certain which are not demonstrated or demonstrable''
\citep[pp.~245--246]{huxley1889}. The claim concerns what can be established, not what it is
reasonable to believe.

\medskip\noindent\textbf{C2. Worldviews that agree on all possible evidence are never
separated.}\phantomsection\label{item:C2}
(From P-M(a) and AA-R2.)
By P-M(a), two worldviews that agree on everything that could ever be experienced are accounts
that agree on all possible evidence. By AA-R2, no further evidence and no ideal judge ever
separates them. The accounts of \emph{Algorithmic Agnosticism} do not speak of God or nature, so applying AA-R2
to worldviews is this paper's extension, made through P-M(a). The hope that future or final
evidence would settle the matter therefore fails for such pairs.

Hick's eschatological verification is a case of that hope. Theism, Hick argued, is verifiable in
principle by experience after death, such as the experience of God's purpose for human beings
fulfilled in communion with God as revealed in Christ, so theism and naturalism are not empirically equivalent
\citep{hick1960}. C2 does not decide whether he was right; it says what follows on each answer. If a
non-theistic account could accommodate every such experience, it and theism agree on all possible
evidence, and C2 applies: the experience, should it come, would not separate them. If not, the two
differ on possible evidence, and C2 does not apply to them. The experience, when it came, would
still be finite evidence with rivals consistent with it, and C1 answers any claim that it
established theism.

\medskip\noindent\textbf{C3. What selects among the worldviews the evidence leaves open is a
choice.}\phantomsection\label{item:C3}
(From P-M(a) and AA-R3.)
By P-M(a), worldviews are accounts, so by AA-R3 a prior probability, simplicity, or any other
criterion that ranks the rivals is an imposed constraint, and the evidence does not determine it.
Probabilistic and parsimony arguments therefore establish nothing beyond the constraint they adopt.

Each decisive element in \cref{sec:arguments} is thus answered: an account or a judgment by C1,
and an imposed constraint, a prior included, by C3.

The conclusions are conditional on P-M, the premise AA-P1 of \emph{Algorithmic Agnosticism}, and logic. D1 and D3
are definitions, and its results are proved there. The conclusions concern
arguments and evidence, not what ultimately is. P-M is not established by D3's standard either;
it is owned as a premise, and the inference from it is logic.

\section{Objections}\label{sec:objections}

The objections are to the premises, the classification, and the scope of the conclusions.

\medskip\noindent\textbf{O1. ``Religious beliefs are not hypotheses.''}\phantomsection\label{item:O1}
On this view it is a conceptual, grammatical truth that the sense of a religious belief lies in
the role it plays in a form of life: in worship, confession, thanksgiving and the way a believer
meets the world. It does not lie in any
explanatory claim about what is. To treat belief in God as a hypothesis that might be confirmed
or disconfirmed is to misread its grammar \citep{phillips1976}.

On this view such beliefs make no claim of the kind an account makes, so P-M(a) does not cover
them and the paper has no quarrel with them. If the thesis is offered as true of what religious
believers actually mean, it is an account that goes beyond the evidence about religious practice,
and by P-M(b) it is not established.

\medskip\noindent\textbf{O2. Basic belief.}\phantomsection\label{item:O2}
On Reformed epistemology, belief in God can be properly basic: formed directly by a cognitive
faculty, not inferred from other beliefs. It is warranted if that faculty functions properly, in
an appropriate environment, according to a design plan aimed at truth. If theism is true, such
belief is probably warranted, and its warrant owes nothing to argument \citep{plantinga2000}.

The paper is silent on rationality, proper function and warrant. By P-M(a), a faculty supplies a
judgment, and ``this faculty is reliable'' is an account, so by P-M(b) it is not established. On externalist theories of knowledge, the belief
may count as knowledge. Whether it does turns on whether the worldview is true, and C1 says that
cannot be established.

\medskip\noindent\textbf{O3. Faith as certain knowledge.}\phantomsection\label{item:O3}
Faith, on this view, is more certain than wisdom, science or understanding in respect of its
cause, because it is founded on the divine truth \citep[II-II q.4 a.8]{aquinas-st}. It is a firm and
certain knowledge of God's benevolence toward us \citep[III.2.7]{calvin1559}. Its assent is given
on the authority of God who reveals, who ``can neither deceive nor be deceived''
\citep[ch.~3]{vatican1870}.

By P-M(a), faith offered as firmness of assent is an imposed constraint: a choice. Faith offered as
knowledge rests on grounds that are accounts, and by P-M(b) they are not established.

\medskip\noindent\textbf{O4. Factive access and unfakeable marks.}\phantomsection\label{item:O4}
Epistemological disjunctivism holds that in the good case a person's reason for belief can be
factive and reflectively available at once: the reason is that one sees that p, not that it
seems to one that p \citep{pritchard2012}. Some traditions hold that certain experiences of God
carry a mark that cannot be counterfeited. Consolation without preceding cause, on Ignatius's
rules, is given by God alone \citep[\S330]{ignatius-exercises}.

By P-M(a) and AA-D, what such an experience supplies is evidence: how things seem. The claim that this seeming is a factive reason, or is
uncounterfeitable, is an account. By P-M(b) it has a consistent rival. Divine veracity is part of
the account, so it cannot establish the account.

\medskip\noindent\textbf{O5. ``God is beyond all categories.''}\phantomsection\label{item:O5}
On apophatic theology, God is beyond every assertion and every denial, so no positive claim about
God is true as stated \citep{dionysius-mt}.

Then the worldview asserts nothing that could be established. Once a positive claim is made, C1
applies.

\medskip\noindent\textbf{O6. ``Only God knows.''}\phantomsection\label{item:O6}
The objection is that God knows which account is true, so ``no one'' is too strong.

Grant that, if God exists, God knows. ``No one'' in the core claim means no one who reasons from
finite evidence.

\medskip\noindent\textbf{O7. ``The standard is too high.''}\phantomsection\label{item:O7}
This is the triviality objection. If establishing a claim means excluding every consistent rival,
no empirical claim is ever established. The core claim then says nothing special about religion.

Concede the parity. By P-M(b) and D3, no account that goes beyond the evidence is established,
including ordinary empirical ones. But
the standard is the one the classic project claims to meet. God, on that project, ``can be known
with certainty'' by the natural light of reason from created things \citep[ch.~2]{vatican1870},
and some critics claim a demonstration of the denial \citep{mackie1955,findlay1948}.

\medskip\noindent\textbf{O8. Atheism as absence of belief.}\phantomsection\label{item:O8}
On this view atheism is not a claim but the lack of one, and so it needs no defense.

Absence of belief makes no claim, so C1 does not apply to it. The denial that God exists is an
account, so C1 applies to it.

\medskip\noindent\textbf{O9. \emph{A priori} premises.}\phantomsection\label{item:O9}
A defender may hold that a premise classed as an imposed constraint is really a truth of reason.
An example is the possibility of a maximally great being, which is rationally acceptable on its
defender's view \citep{plantinga1974}.

If it were a truth of reason, and it excluded every rival, P-M(b) would be false for that
worldview, and C1 would not reach it.

\medskip\noindent\textbf{O10. ``The mapping begs the question.''}\phantomsection\label{item:O10}
The objection is that P-M(b) assumes what the paper sets out to show.

P-M(b) is an owned premise, with its refutation condition (\cref{item:PM}, with its motivation). The classic \emph{a priori} arguments are
attempts to meet its refutation condition, and \cref{sec:arguments} names, for each argument, the
element it would have to show to be a truth of reason (O9).

\section{Conclusion}\label{sec:conclusion}

No argument from experience or reason establishes that any worldview, whether theism, naturalism
or the skeptical hypothesis, is the true account of what is. Any of them may be true, but no one
can establish that it is. No evidence ever separates worldviews that agree on all possible
evidence, and what selects among the worldviews the evidence leaves open is a choice.

The core claim is not that God does not exist; it concerns what can be established, not what it is
reasonable to believe.

Hawking hoped that a complete theory would let us know the mind of God \citep{hawking1988}. Even a
complete theory would be an account. We might be right about the mind of God; we could
never establish that we were.

\section*{Statement on the use of AI}
\addcontentsline{toc}{section}{Statement on the use of AI}

This paper was written with substantial assistance from AI systems: Claude Code
(Anthropic; Claude Opus~5.5) and Grok (xAI). They were used to develop and review the argument,
to research and check the literature, and to draft the text. AI systems are not authors. The
author directed the work and is responsible for its content. A record of the process is in the
project repository.



\bibliographystyle{plainnat}
\bibliography{refs}

\end{document}
